\scanpage{33}
\mspage{\unsure{37|31}}
En effet, on peut supposer $k$ algébriquement clos, de sorte que $T$ est contenu dans un tore maximal\deleted{\unsure{max}} $T'$. Soit $C'$ le centralisateur de ce dernier, alors on a $C'\subset C$, donc $\dim C'\leq\dim C$, cqfd.

b) \deleted{Le raisonnement de a) nous montre que} Si $\rho_r$ est localement constant, alors pour tout tore $T$ dans $G$, \deleted{et} et tout $s\in S$, si $T_s$ est \inserted{tore} maximal dans $G_s$, alors il existe un voisinage ouvert $U$ de $s$ tel que $T|U$ soit tore maximal dans $G\ U$. Utilisant maintenant le raisonnement de a), on voit que \deleted{\unsure{si}} (i)$\Longrightarrow$(ii bis). D'autre part (ii bis)$\Longrightarrow$(i), car \deleted{on peut supposer $S$} lorsque $G$ admet un tore maximal $T$, alors il est évident que $\rho_r(s)=\dim T_s$ est une fonction localement constante de \deleted{\unsure{s}} $s$, or nous avons signalé au début de la démonstration de a) que la question de la continuité de $\rho_r$ était locale pour la topologie fidèlement plate. Donc (i)$\Longleftrightarrow$(ii bis), et évidemment (ii)$\Longrightarrow$(ii bis), reste à prouver l'implication inverse \inserted{(i)$\Longrightarrow$(ii)}. Pour ceci, introduisons le foncteur $F$ de \inserted{Exp $\overline{\underline{\mathrm{XI}}}$} 4.1., qui est donc un préschéma \deleted{noethérien} lisse \inserted{et séparé} sur $S$, et considérons le sous-foncteur $\unsure{\mathcal T|\tau}$ de $F$, dont la valeur \deleted{sur} pour un $S'$ sur $S$ est l'ensemble des tores maximaux dans $G_{S'}$. Utilisant la remarque précédente \deleted{que}\inserted{que, moyennant (i),} un tore dans $G_{S'}$ qui est maximal dans la fibre d'un point $s\in S'$ est maximal au dessus d'un voisinage ouvert de $s$, on voit que $\unsure{\mathcal T|\tau}$ est représentable par un sous-préschéma ouvert de $F$, et par suite est lisse \inserted{et séparé} sur $S$. Comme le morphisme structural $\unsure{\mathcal T|\tau}\longrightarrow S$ est évidemment surjectif, il admet donc \deleted{localement} une section localement au sens de la topologie étale en vertu de \inserted{Exp $\overline{\underline{\mathrm{XI}}}$} 1.10., ce qui prouve (i)$\Rightarrow$(ii).
\margin{$\unsure{\mathcal T|\tau}$}

c) C'est une conséquence immédiate de \inserted{Exp $\overline{\underline{\mathrm{XI}}}$} 5.4. bis, compte tenu du théorème de conjugaison de Borel rappelé au début du numéro.

d) Moyennant la remarque du début de la démonstration dans a), et compte tenu de b), on peut supposer qu'il existe un tore maximal $T$
\mspage{32}
